\documentclass[11pt, a4paper]{article}

% ── PACKAGES ──────────────────────────────────────────────────
\usepackage[utf8]{inputenc}
\usepackage[T1]{fontenc}
\usepackage{geometry}
\geometry{
  top=2.5cm,
  bottom=2.5cm,
  left=2.5cm,
  right=2.5cm
}
\usepackage{hyperref}
\hypersetup{
  colorlinks=true,
  urlcolor=blue,
  linkcolor=blue,
  citecolor=blue,
  pdftitle={Mars Biosphere Collapse vs. Emergence:
    A Pre-Registered Geometric Prediction of the
    Complexity Gradient in the Mars Sedimentary Record
    Prior to Sample Return},
  pdfauthor={Eric Robert Lawson}
}
\usepackage{microtype}
\usepackage{parskip}
\usepackage{booktabs}
\usepackage{array}
\usepackage{tabularx}
\usepackage{xcolor}
\usepackage{mdframed}
\usepackage{amsmath}
\usepackage{amssymb}
\usepackage{enumitem}
\usepackage{titlesec}
\usepackage{lmodern}
\usepackage{authblk}
\usepackage{abstract}

% ── COLOURS ───────────────────────────────────────────────────
\definecolor{headerblue}{RGB}{20,60,140}
\definecolor{boxblue}{RGB}{20,60,140}
\definecolor{rulecolor}{RGB}{20,60,140}

% ── SECTION FORMAT ────────────────────────────────────────────
\titleformat{\section}
  {\large\bfseries\color{headerblue}}
  {\thesection.}{0.5em}{}
\titleformat{\subsection}
  {\normalsize\bfseries\color{headerblue}}
  {\thesubsection.}{0.5em}{}
\titleformat{\subsubsection}
  {\normalsize\itshape\color{headerblue}}
  {\thesubsubsection.}{0.5em}{}

\setlength{\parskip}{0.6em}
\setlength{\parindent}{0pt}

% ══════════════════════════════════════════════════════════════
\begin{document}

% ── TITLE BLOCK ───────────────────────────────────────────────
\title{
  \large\bfseries\color{headerblue}
  MARS BIOSPHERE COLLAPSE VS.\ EMERGENCE\\[0.4em]
  \normalsize\color{headerblue}
  A Pre-Registered Geometric Prediction of the
  Complexity Gradient\\
  in the Mars Sedimentary Record Prior to Sample
  Return,\\
  the Water Attractor State Cascade Model,\\
  and the Two-Axis LECA Stabilisation Framework
}

\author[1]{Eric Robert Lawson}
\affil[1]{OrganismCore (Independent Research)\\
  ORCID: \href{https://orcid.org/0009-0002-0414-6544}
  {0009-0002-0414-6544}\\
  \href{mailto:OrganismCore@proton.me}
  {OrganismCore@proton.me}
}

\date{2026-03-13\\[0.3em]
  \footnotesize
  LECA Pre-registration DOI:
  \href{https://doi.org/10.5281/zenodo.18986790}
  {10.5281/zenodo.18986790}
  (timestamp: 2026-03-12,
  prior to this derivation)\\
  Universal Life Detection Standard DOI:
  \href{https://doi.org/10.5281/zenodo.18994082}
  {10.5281/zenodo.18994082}\\
  Repository:
  \href{https://github.com/Eric-Robert-Lawson/attractor-oncology}
  {github.com/Eric-Robert-Lawson/attractor-oncology}
}

\maketitle
\thispagestyle{empty}

% ── STATEMENT BOX ─────────────────────────────────────────────
\begin{mdframed}[
  backgroundcolor=boxblue,
  linecolor=boxblue,
  linewidth=0pt,
  innertopmargin=0.5em,
  innerbottommargin=0.5em,
  innerleftmargin=1em,
  innerrightmargin=1em
]
\begin{center}
\small\bfseries\color{white}
The Cheyava Falls fossil record (Hurowitz et al.,
\textit{Nature} 645, 2025) is consistent with either
the beginning or the end of a Martian biosphere.
These two interpretations make opposite predictions
about the complexity gradient in the Mars subsurface
sedimentary record: if Cheyava Falls represents
biosphere \textit{emergence}, complexity increases
with decreasing depth; if it represents biosphere
\textit{collapse}, complexity increases with
increasing depth and age.
This prediction is derived from geometric first
principles, is locked here before Mars sample return,
and is testable by the Mars sample return mission
(2030s target).
Pre-registered: 2026-03-13.
\end{center}
\end{mdframed}

\vspace{1em}

% ── ABSTRACT ──────────────────────────────────────────────────
\begin{abstract}
\noindent
We present a pre-registered geometric prediction
concerning the directional ambiguity in the Cheyava
Falls biosignature record on Mars.
The magnetic field loss of Mars ($\sim$4.0--4.1\,Ga)
and the Cheyava Falls deposition age
($\sim$3.8--4.1\,Ga) are coincident within
measurement error, creating a fundamental directional
ambiguity: the fossil may represent life establishing
itself (emergence) or life in terminal decline
(collapse).
We derive two mutually exclusive predictions from
attractor geometry that are testable by the complexity
gradient in the Mars subsurface sedimentary record,
to be resolved by Mars sample return.
We additionally present the water attractor state
cascade model, which identifies three distinct
chemical attractor states of planetary water chemistry
and the biological grades each supports.
We present the two-axis LECA stabilisation model,
which identifies independent nitrogen and oxygen
destabilisation axes that together define the
conditions under which the LECA attractor is the
stable state of carbon chemistry.
All derivations are from geometric first principles.
This document is the pre-registration of the
complexity gradient prediction before the Mars sample
return mission retrieves material from the Jezero
Crater sedimentary record.
\end{abstract}

\vspace{0.5em}
\noindent\textcolor{rulecolor}{\rule{\linewidth}{0.6pt}}

% ── SECTION 1 ─────────────────────────────────────────────────
\section{The Timestamp Collision}

Two facts are established independently of this
framework:

\begin{enumerate}[leftmargin=*, label=(\arabic*)]
\item The Mars magnetosphere was lost approximately
  4.0--4.1\,Ga, as determined by MAVEN mission crustal
  magnetisation measurements.
\item The Cheyava Falls biosignature features in the
  Bright Angel formation of Jezero Crater are dated
  approximately 3.8--4.1\,Ga, as reported in
  Hurowitz et al.\ \cite{hurowitz2025}.
\end{enumerate}

These ranges overlap within measurement error.
The magnetic field loss and the depositional age of
the fossil features occurred within the same
geological event window.

This creates a directional ambiguity that cannot be
resolved from surface data alone.

The fossil record has no direction arrow.
It records what was present.
It does not record whether what was present was
growing or dying.

% ── SECTION 2 ─────────────────────────────────────────────────
\section{The Two Interpretations}

\subsection{Option A --- Emergence}

Under this interpretation, the Cheyava Falls lake
represents a newly habitable environment.
Life was forming or recently established in
State~1 water conditions.
The LUCA attractor was establishing in a new
depositional environment.
The fossil is the beginning of a cascade that
was then interrupted by atmospheric collapse.

\textbf{Complexity gradient prediction (Option A):}
Biosignature complexity \textit{decreases} with
increasing depth and age.
Simpler organisms deeper and earlier.
More complex features shallower and later.
Cheyava Falls is the most developed record
in the sequence.

\subsection{Option B --- Collapse}

Under this interpretation, Mars harboured a more
complex biosphere that predated Cheyava Falls.
The magnetic field loss initiated atmospheric
stripping. Conditions deteriorated progressively.
Complex organisms, requiring the most stable
conditions, died first.
Simpler chemotrophic organisms survived longest.
LUCA-grade iron-sulfur chemotrophs were the last
surviving life form.
The Cheyava Falls leopard spots are the final
record of a dying world.

\textbf{Complexity gradient prediction (Option B):}
Biosignature complexity \textit{increases} with
increasing depth and age.
More complex features deeper and earlier.
Simpler features shallower and later.
Cheyava Falls is the least developed record
in the sequence --- the last survivor.

\subsection{Why Current Data Cannot Decide}

The surface record is the survival record,
not the complete record.

Three billion years of unshielded solar and
cosmic ray bombardment (post-magnetosphere loss
at 4.0--4.1\,Ga) preferentially destroys complex
organic biosignatures.
Simple iron-sulfur minerals --- vivianite,
greigite --- survive.
Complex organic molecules, proteins, and nucleic
acids do not.

Whatever complexity existed at Cheyava Falls
at time of deposition has been partially or
entirely degraded by radiation.

The surface record is the floor of complexity,
not the ceiling.

Additionally, Mars never underwent a
Great Oxidation Event (GOE) equivalent.
The entire water column remained anoxic
throughout the habitable period.
This produces colony-scale mineral reaction
fronts as the dominant fossilisation mechanism
rather than individually resolved cellular
fossils (see Section~5).

Both Option~A and Option~B are consistent
with the currently available surface data.

% ── SECTION 3 ─────────────────────────────────────────────────
\section{The Pre-Registered Prediction}

\begin{mdframed}[
  linecolor=rulecolor,
  linewidth=0.8pt,
  innertopmargin=0.6em,
  innerbottommargin=0.6em,
  innerleftmargin=1em,
  innerrightmargin=1em
]
\textbf{Pre-Registered Prediction (locked 2026-03-13,
prior to Mars sample return):}

\medskip
The complexity gradient in the Jezero Crater
sedimentary record with depth and increasing age
will run in one of two directions:

\medskip
\textbf{If collapse (Option B):}
Subsurface samples at greater stratigraphic depth
than the Cheyava Falls surface sample will show
\textit{more} complex biosignatures, specifically:
\begin{itemize}[leftmargin=*]
\item More diverse mineral assemblages.
\item Richer and more varied organic carbon content.
\item More complex spatial organisation of
  mineral--organic associations.
\item Biosignature features at smaller scale
  (approaching individual cell scale of 10--30\,$\mu$m
  rather than colony scale of 1--2\,mm)
  if preservation conditions permit.
\end{itemize}

\textbf{If emergence (Option A):}
Subsurface samples at greater stratigraphic depth
will show \textit{less} complex biosignatures,
consistent with earlier, simpler biological
communities.

\medskip
\textbf{The Mars sample return mission resolves
which prediction is correct.}
\end{mdframed}

% ── SECTION 4 ─────────────────────────────────────────────────
\section{The Water Attractor State Cascade Model}

\subsection{Water Chemistry as Attractor State}

Water is not a uniform medium.
The chemical composition of a body of water
constitutes an attractor state of the water
chemistry itself.
Different attractor states support different
grades of biological complexity.
The cascade from first life to a biologically
maintained atmosphere is a cascade of coupled
water attractor states.

\subsection{The Three States}

\begin{table}[h]
\centering
\small
\begin{tabularx}{\linewidth}{p{1.2cm}p{3.0cm}p{3.0cm}X}
\toprule
\textbf{State} & \textbf{Chemistry} &
\textbf{Life supported} & \textbf{Transition} \\
\midrule
State 1 &
pH 5--6; high Fe$^{2+}$; anoxic;
low N; pre-GOE &
LUCA-grade chemolithotrophs only.
Cannot support stable LECA. &
Biological photosynthesis
drives GOE. \\[0.3em]
State 2 &
pH 6.5--7.5; Fe$^{2+}$ declining;
O$_2$ rising; N cycling increasing &
LECA-grade eukaryotes.
Mitochondria functional.
LECA attractor stable. &
LECA diversification drives
State~3. \\[0.3em]
State 3 &
pH 8.1; trace Fe; high sulfate;
O$_2$ abundant &
Full triadic cycle.
Biologically maintained atmosphere. &
Terminal state on Earth. \\
\bottomrule
\end{tabularx}
\caption{The three water attractor states on Earth.
  Each transition is driven by biological activity
  in the previous state.}
\end{table}

\subsection{Mars in This Framework}

The Cheyava Falls mineralogy (vivianite =
Fe$_3$(PO$_4$)$_2$; greigite = Fe$_3$S$_4$;
anoxic conditions; organic carbon) confirms
State~1 water chemistry in the Jezero Crater
lake environment \cite{hurowitz2025}.

Mars had State~1 water.
The magnetic field loss at 4.0--4.1\,Ga
prevented oxygen accumulation.
No GOE equivalent occurred.
The cascade was interrupted at State~1.
Mars never reached State~2.

% ── SECTION 5 ─────────────────────────────────────────────────
\section{The Two-Axis LECA Stabilisation Model}

\subsection{Two Independent Destabilisation Axes}

The LECA attractor is stabilised by the
\textit{simultaneous} absence of two independent
destabilisation signals.

\textbf{Axis 1 --- Nitrogen:}
Bioavailable nitrogen (NH$_4^+$, NO$_3^-$, or
amino acids above threshold) fires germination
and multicellular organisation programs.
In the LECA arrest medium: explicitly absent.
Tested in ARM~A by nitrogen release control at
72 hours (addition of (NH$_4$)$_2$SO$_4$
restores germination).

\textbf{Axis 2 --- Oxygen:}
Sufficient O$_2$ (above $\sim$3\% for microaerophilic
threshold) activates aerobic metabolic programs
unavailable in pre-GOE conditions.
Mitochondria at full aerobic capacity provide
the energy density to advance developmental
programs past $D{=}0$.
In the LECA arrest medium: held at $<$3\% O$_2$
throughout.

\subsection{What the GOE Did}

On Earth, the GOE ($\sim$2.4\,Ga) raised Axis~2
above the destabilisation threshold.
Concurrent nitrogen cycling changes raised Axis~1.
Both axes exceeded their thresholds together.
The LECA attractor was destabilised at
planetary scale.
Complex eukaryotic life emerged.

On Mars, the magnetic field loss prevented
oxygen accumulation.
Axis~2 never rose above threshold.
Nitrogen cycling could not change.
Axis~1 remained below threshold.
The LECA attractor remained stable in an
increasingly hostile and ultimately uninhabitable
planetary environment.

\subsection{Impact on the Arrest Medium}

This analysis confirms that the LECA Resurrection
Protocol arrest medium (pre-registered DOI:
10.5281/zenodo.18986790) correctly addresses both
axes.
No change to the protocol is required.
The two-axis understanding deepens the
mechanistic explanation of why the conditions are
specified as they are.

% ── SECTION 6 ─────────────────────────────────────────────────
\section{The Fossilisation Paradox and Its Resolution}

On Earth, the best cellular fossils are preserved
at redox boundaries between an oxidising surface
layer and anoxic deep sediment.
This produces sharp mineralisation of organic
material at the boundary, sometimes preserving
cellular detail.

Mars never had a GOE equivalent.
The entire water column remained anoxic.
No sharp oxidising--anoxic boundary existed.

The predicted fossilisation mechanism for
State~1 water conditions is therefore:

\begin{itemize}[leftmargin=*]
\item Not individual cellular-detail fossils.
\item Colony-scale mineral reaction fronts.
\item The bulk mineral product of microbial
  community metabolism at population scale.
\item Vivianite and greigite formed by
  Fe-phosphate and Fe-sulfide reducing metabolism
  at 1--10\,mm scale (colony scale) rather than
  10--30\,$\mu$m scale (individual cell scale).
\end{itemize}

This is not a failure of the biological interpretation
of Cheyava Falls.
It is the geometric prediction of what biosignatures
look like on a world that never underwent a GOE.

The leopard spots (1--2\,mm, vivianite--greigite)
are consistent with this mechanism.
The poppy seeds ($<$200\,$\mu$m) approach colonial
microstructure scale.
Individual cellular-scale confirmation (10--30\,$\mu$m)
requires Mars sample return resolution.

% ── SECTION 7 ─────────────────────────────────────────────────
\section{The ARM A Reference Standard for Sample
Return}

The LECA Resurrection Protocol ARM~A
(pre-registered DOI: 10.5281/zenodo.18986790,
timestamp: 2026-03-12) will produce a formally
defined, pre-registered morphological profile of
a living LECA-grade cell.

This profile constitutes the formal reference
standard against which Mars sample return material
should be compared.

The ARM~A reference standard:
\begin{itemize}[leftmargin=*]
\item Is pre-registered before the Mars samples
  arrive on Earth.
\item Is derived from the same geometric framework
  that generated the Cheyava Falls prediction.
\item Is physically anchored by laboratory
  confirmation in living biological material.
\item Formally defines the expected cellular
  diameter (10--30\,$\mu$m), nuclear morphology
  (DAPI-stained nucleus confirmed), and validation
  criteria for LECA-grade cell identification.
\end{itemize}

If the Cheyava Falls sample return material
contains sub-features at 10--30\,$\mu$m scale
matching the ARM~A morphological profile, that
constitutes individual cell-scale confirmation
of LECA-grade life on ancient Mars.

The relationship between ARM~A ($\sim$\$200,
BSL-1, days to result) and Mars sample return
(multi-billion dollar mission, 2030s) is explicit:
the cheap experiment in Rogers Park, Chicago
provides the calibration standard for the most
consequential planetary science mission in
human history.

% ── SECTION 8 ─────────────────────────────────────────────────
\section{Summary of Pre-Registered Predictions}

\begin{table}[h]
\centering
\small
\begin{tabularx}{\linewidth}{p{0.5cm}p{4.5cm}p{2.0cm}X}
\toprule
\textbf{\#} & \textbf{Prediction} &
\textbf{Test} & \textbf{Resolves} \\
\midrule
P1 &
Complexity gradient increases with depth
if collapse (Option B). &
Mars sample return.
Jezero Crater subsurface. &
Whether Cheyava Falls is a death
record or a birth record. \\[0.3em]
P2 &
Complexity gradient decreases with depth
if emergence (Option A). &
Mars sample return. &
Same. \\[0.3em]
P3 &
Sub-features at 10--30\,$\mu$m in the
poppy seeds matching the ARM~A profile
confirm individual LECA-grade cells. &
Mars sample return
nanometre-resolution
imaging. &
Individual cell-scale
confirmation of LECA-grade
life on ancient Mars. \\
\bottomrule
\end{tabularx}
\caption{Pre-registered predictions locked 2026-03-13,
  prior to Mars sample return.}
\end{table}

% ── SECTION 9 ─────────────────────────────────────────────────
\section{Honest Statement of Limitations}

\begin{enumerate}[leftmargin=*, label=(\arabic*)]
\item \textbf{Surface data is insufficient.}
  The collapse vs.\ emergence distinction cannot
  be resolved from Perseverance rover surface data.
  It requires access to the stratigraphic depth
  record in the Cheyava Falls formation.
  Mars sample return provides this.

\item \textbf{Radiation taphonomy is uncertain.}
  3.8 billion years of unshielded radiation have
  preferentially degraded complex biosignatures.
  Even under Option~B (collapse), the subsurface
  record may not preserve sufficient complexity
  contrast to distinguish the two interpretations.
  Radiation shielding at depth is the critical
  variable.

\item \textbf{ARM~A pending.}
  The morphological reference standard requires
  ARM~A confirmation.
  Until ARM~A runs, the reference standard is
  specified but not physically anchored.

\item \textbf{Common origin vs.\ independent
  emergence unresolved.}
  This document does not address whether any
  Martian life shared common ancestry with Earth
  life. Both scenarios are consistent with the
  geometric prediction. The distinction requires
  genomic analysis of Mars sample return material.
\end{enumerate}

% ── SECTION 10 ────────────────────────��───────────────────────
\section{Conclusion}

The Cheyava Falls fossil record is ambiguous in
direction. The ambiguity is geometric, not
observational, and cannot be resolved from
currently available data.

Two mutually exclusive predictions follow from
attractor geometry:

\begin{itemize}[leftmargin=*]
\item Under biosphere collapse (Option~B):
  complexity increases with depth.
\item Under biosphere emergence (Option~A):
  complexity decreases with depth.
\end{itemize}

This prediction is locked here, on 2026-03-13,
before the Mars sample return mission retrieves
subsurface material from Jezero Crater.

The water attractor state cascade model and the
two-axis LECA stabilisation model are presented
as the geometric basis for understanding why
State~1 water conditions on ancient Mars were
the conditions in which the LUCA attractor was
the stable biological state.

ARM~A provides the morphological reference
standard that makes the Mars sample return
comparison formally defined rather than
interpretively open.

The geometry was derived before the data.

The data will confirm which geometric prediction
is correct.

Both outcomes are informative.

Neither can be adjusted after the fact.

This document is the lock.

\vspace{1em}
\noindent\textcolor{rulecolor}{\rule{\linewidth}{0.4pt}}

% ── ACKNOWLEDGMENTS ───────────────────────────────────────────
\section*{Acknowledgments}

This work was conducted independently without
institutional funding or affiliation.
The geometric derivation sequence is documented
in the pre-registration record
(DOI: 10.5281/zenodo.18986790) and the
attractor-oncology repository
(\href{https://github.com/Eric-Robert-Lawson/attractor-oncology}{github.com/Eric-Robert-Lawson/attractor-oncology}).
The Cheyava Falls data is the work of the NASA
Perseverance mission science team
(Hurowitz et al., \textit{Nature}, 2025)
and is referenced here as independent
contextual data, not as a derivation source.

% ── REFERENCES ────────────────────────────────────────────────
\begin{thebibliography}{99}

\bibitem{preregistration}
Lawson, E.R.\ (2026).
\textit{LECA Resurrection Protocol ---
Pre-Registration v1.0}.
Zenodo.
\href{https://doi.org/10.5281/zenodo.18986790}
{doi:10.5281/zenodo.18986790}

\bibitem{amendmentA1}
Lawson, E.R.\ (2026).
\textit{LECA Resurrection Protocol ---
Registered Amendment A1}.
Zenodo.
\href{https://doi.org/10.5281/zenodo.18989573}
{doi:10.5281/zenodo.18989573}

\bibitem{lucapaper}
Lawson, E.R.\ (2026).
\textit{The Ribosome Is the LUCA}.
Zenodo.
\href{https://doi.org/10.5281/zenodo.18988844}
{doi:10.5281/zenodo.18988844}

\bibitem{uldstandard}
Lawson, E.R.\ (2026).
\textit{The Universal Life Detection Standard}.
Zenodo.
\href{https://doi.org/10.5281/zenodo.18994082}
{doi:10.5281/zenodo.18994082}

\bibitem{hurowitz2025}
Hurowitz, J.A., Tice, M.M., Allwood, A.C.,
et al.\ (2025).
Redox-driven mineral and organic associations
in Jezero Crater, Mars.
\textit{Nature}, 645(8080), 332--340.
\href{https://doi.org/10.1038/s41586-025-09413-0}
{doi:10.1038/s41586-025-09413-0}

\bibitem{margulis1967}
Sagan, L.\ (1967).
On the origin of mitosing cells.
\textit{Journal of Theoretical Biology},
14(3), 225--274.
\href{https://doi.org/10.1016/0022-5193(67)90079-3}
{doi:10.1016/0022-5193(67)90079-3}

\bibitem{nealson1997}
Nealson, K.H.\ \& Stahl, D.A.\ (1997).
Microorganisms and biogeochemical cycles.
In: Banfield, J.F.\ \& Nealson, K.H.\ (Eds.),
\textit{Geomicrobiology: Interactions between
Microbes and Minerals}.
Reviews in Mineralogy, Vol.\ 35, pp.\ 5--34.
Mineralogical Society of America.

\bibitem{lane2015}
Lane, N.\ (2015).
\textit{The Vital Question}.
W.W.\ Norton.
ISBN: 978-0-393-08881-6

\end{thebibliography}

% ── FOOTER ────────────────────────────────────────────────────
\vspace{1em}
\noindent\textcolor{rulecolor}{\rule{\linewidth}{0.4pt}}
\begin{center}
\footnotesize
\textit{The geometry came first.\ \
The data will say which prediction is correct.\ \
Both outcomes are informative.\ \
Neither can be adjusted.}\\[0.15em]
\href{https://github.com/Eric-Robert-Lawson/attractor-oncology}
{\texttt{github.com/Eric-Robert-Lawson/attractor-oncology}}
\ $\cdot$\
\href{https://doi.org/10.5281/zenodo.18986790}
{\texttt{doi.org/10.5281/zenodo.18986790}}
\ $\cdot$\
\href{https://orcid.org/0009-0002-0414-6544}
{\texttt{ORCID: 0009-0002-0414-6544}}
\end{center}

\end{document}
